\documentclass[11pt]{article}
\usepackage[margin=1in]{geometry}
\usepackage{amsmath,amssymb,amsthm}
\usepackage[hidelinks]{hyperref}

\newtheorem{theorem}{Theorem}
\newtheorem{lemma}{Lemma}
\newtheorem{remark}{Remark}

\newcommand{\C}{\mathbb C}
\newcommand{\Z}{\mathcal Z}
\newcommand{\otg}{\otimes_\gamma}
\newcommand{\Ann}{\operatorname{Ann}}

\title{A Finite-Dimensional Counterexample to Tensoring Centers for Banach Algebras}
\author{FA Banach run \texttt{fa\_banach\_001}, agent \texttt{agent\_lane\_19}}
\date{2026-06-20}

\begin{document}
\maketitle

\section*{Status}

This packet records a finite-dimensional counterexample to the unrestricted
Banach-algebra analogue of the center theorem proved by Gupta and Jain for
C*-algebras. The source paper proves that, for C*-algebras \(A\) and \(B\),
\[
  \Z(A\otg B)=\Z(A)\otg \Z(B),
\]
and then asks whether there is a similar relationship for Banach algebras.

The example below shows that the direct equality fails for Banach algebras, even
in dimension \(4\cdot 2\).

\section*{Source Question}

After proving the C*-algebra identity
\[
  \Z(A\otg B)=\Z(A)\otg\Z(B),
\]
Gupta and Jain ask:
\begin{quote}
Is there any relationship, as above, between the center of \(A\otg B\) and
\(\Z(A)\otg\Z(B)\) for Banach algebras \(A\) and \(B\)?
\end{quote}

\section*{A General Annihilator Mechanism}

For a Banach algebra \(B\), write
\[
  \Ann(B)=\{b\in B: bB=Bb=\{0\}\}.
\]

\begin{lemma}
Let \(A\) and \(B\) be Banach algebras. If \(0\neq b\in\Ann(B)\), then
\[
  A\otimes \C b \subseteq \Z(A\otg B).
\]
Consequently, if \(A\) is noncommutative and \(b\neq0\), then the center of
\(A\otg B\) may contain elements not coming from \(\Z(A)\otg\Z(B)\).
\end{lemma}

\begin{proof}
For elementary tensors,
\[
  (a\otimes b)(x\otimes y)=ax\otimes by=0,\qquad
  (x\otimes y)(a\otimes b)=xa\otimes yb=0
\]
for every \(a,x\in A\) and \(y\in B\). Hence every \(a\otimes b\) commutes with
every elementary tensor. Since multiplication on \(A\otg B\) is continuous and
the algebraic tensor product is dense, \(a\otimes b\) is central in
\(A\otg B\).
\end{proof}

\section*{The Counterexample}

\begin{theorem}
There are finite-dimensional Banach algebras \(A\) and \(B\) such that
\[
  \Z(A)\otg\Z(B) \subsetneq \Z(A\otg B).
\]
\end{theorem}

\begin{proof}
Let \(A=M_2(\C)\) with its usual operator norm. Let \(B=\C e\oplus \C n\) with
multiplication
\[
  e^2=e,\qquad en=ne=n^2=0,
\]
and norm \(\|\alpha e+\beta n\|=|\alpha|+|\beta|\). This is a
finite-dimensional Banach algebra. It is commutative, so \(\Z(B)=B\), and
\[
  \Z(A)=\C I.
\]

Because all spaces are finite-dimensional, the projective tensor product is the
algebraic tensor product equipped with a complete norm. Every element of
\(A\otg B\) has a unique form
\[
  X\otimes e+Y\otimes n,\qquad X,Y\in A.
\]
Let \(T,S\in A\). The multiplication rules in \(B\) give
\[
  (X\otimes e+Y\otimes n)(T\otimes e+S\otimes n)=XT\otimes e
\]
and
\[
  (T\otimes e+S\otimes n)(X\otimes e+Y\otimes n)=TX\otimes e.
\]
Therefore \(X\otimes e+Y\otimes n\) is central in \(A\otg B\) if and only if
\[
  XT=TX\qquad\text{for every }T\in M_2(\C).
\]
This is equivalent to \(X\in\C I\), while \(Y\in M_2(\C)\) is arbitrary. Hence
\[
  \Z(A\otg B)=\C I\otimes e\;+\;M_2(\C)\otimes n.
\]
On the other hand,
\[
  \Z(A)\otg\Z(B)=\C I\otg B
  =\C I\otimes e\;+\;\C I\otimes n.
\]
The inclusion is strict. For instance, if \(E_{12}\) is the standard matrix
unit, then \(E_{12}\otimes n\in\Z(A\otg B)\). It cannot belong to
\(\Z(A)\otg\Z(B)\): applying the coordinate functional \(e^*\) on \(B\) forces
the \(I\otimes e\) coefficient to be zero, and applying \(n^*\) would then force
\(E_{12}\in\C I\), which is impossible.
\end{proof}

\section*{What This Does and Does Not Settle}

The example gives a negative answer to the direct unrestricted Banach-algebra
extension of the C*-algebra identity. The obstruction is not analytic subtlety
of the projective norm; it is algebraic annihilator mass in a nonunital Banach
algebra.

It does not address stronger formulations in which both Banach algebras are
unital, or more generally have bounded approximate identities. In those cases a
nonzero two-sided annihilator is impossible, so this particular obstruction is
removed.

\section*{Novelty Check}

The local run indexes were searched for arXiv:1809.01131, the source title, and
keywords around centers of projective tensor products of Banach algebras. No
prior packet for this question was present. External searches for the center
question and for the annihilator mechanism found the source paper and adjacent
C*-algebra tensor-product literature, but no direct recorded use of this
finite-dimensional example as an answer to the final Banach-algebra question.

\begin{thebibliography}{9}

\bibitem{GuptaJain2018}
V. P. Gupta and R. Jain,
\emph{On Banach space projective tensor product of C*-algebras},
arXiv:1809.01131, 2018.

\end{thebibliography}

\end{document}
